\item \model{MiniMax-M2}

\paragraph*{رویکرد فعال‌سازی مینیاتوری در تقابل با فشرده‌سازی کوانتیزه}
راهبرد بنیادی حاکم بر سری مدل‌های \model{MiniMax-M2} \cite{minimax2026m2} مبتنی بر این اصل است که به‌جای تکیه بر روش‌های تهاجمی کوانتیزاسیون زیر ۸ بیت (که منجر به افت عملکرد در وظایف چندمرحله‌ای عامل‌ها می‌شود)، باید از تقلیل پارامترهای فعال (\en{Mini Activations}) بهره گرفت. مدل پرچمدار \model{M2} با دارا بودن ۲۲۹.۹ میلیارد پارامتر کل، تنها ۹.۸ میلیارد پارامتر فعال را در هر گام پردازش می‌کند.

\paragraph*{محدودیت‌های کوانتیزیشن در لایه‌های توجه خطی}
گزارش فنی \model{M2} نشان می‌دهد که تلاش‌های پیشین برای اعمال فرمت‌های کم‌دقت بر لایه‌های توجه خطی (نظیر \en{Lightning Attention}) با چالش‌های اساسی زیر مواجه شده است:
\begin{itemize}
    \item \textbf{انباشت پویای خطا در حالت بازگشتی:} حالت بازگشتی توجه خطی $S_t = S_{t-1} + K_t^\top V_t$ در مواجهه با تقریب کوانتیزه به شکل زیر عمل می‌کند:
    \begin{equation}
        \tilde{S}_t = S_t + \sum_{\tau=1}^{t} \mathcal{E}_\tau^{\text{quant}}
    \end{equation}
    این خطای تجمعی به سرعت سیگنال بازیابی اطلاعات در متون طولانی را مخدوش می‌سازد.
    \item \textbf{ناسازگاری با حافظه پنهان پیشوندی:} فقدان ساختار پیشرفته کوانتیزاسیون برای کش‌های هیبریدی، استقرار این مدل‌ها را دشوار می‌کرد.
\end{itemize}
در نتیجه این تحلیل‌ها، مدل \model{M2} معماری توجه را به توجه کامل چندسر با اشتراک سر کلید-مقدار (\en{GQA}) در دقت اعشاری کامل بازگرداند و هزینه حافظه را از طریق الگوریتم دقیق ادغام درخت پیشوند (\en{Prefix Tree Merging}) با \textbf{صفر درصد خطای تقریب} تا سقف $40\times$ کاهش داد.

\begin{figure}[htbp]
\centering
\resizebox{0.9\textwidth}{!}{%
\begin{tikzpicture}[
    node distance=1.2cm and 1.6cm,
    ptree/.style={rectangle, draw=violet!70!black, fill=violet!5, rounded corners, minimum width=3.4cm, minimum height=1.1cm, align=center, font=\small\bfseries},
    seq/.style={rectangle, draw=teal!70!black, fill=teal!10, rounded corners, minimum width=2.6cm, minimum height=1.0cm, align=center, font=\small},
    arr/.style={-Latex, thick, draw=violet!80!black}
]
    \node[ptree] (prefix) {پیشوند مشترک طولانی\\ \en{Common Context Prefix}};
    \node[seq, above right=0.3cm and 1.5cm of prefix] (seq1) {شاخه پاسخ ۱};
    \node[seq, right=1.5cm of prefix] (seq2) {شاخه پاسخ ۲};
    \node[seq, below right=0.3cm and 1.5cm of prefix] (seq3) {شاخه پاسخ ۳};

    \draw[arr] (prefix) -- (seq1);
    \draw[arr] (prefix) -- (seq2);
    \draw[arr] (prefix) -- (seq3);
\end{tikzpicture}%
}
\caption{الگوی ادغام درخت پیشوند (\en{Prefix Tree Merging}) به عنوان جایگزین بدون تلفات کوانتیزاسیون کش در \model{MiniMax-M2}}
\label{fig:minimax_m2_tree}
\end{figure}